\title{DeepSeekMath: Pushing the Limits of Mathematical Reasoning in Open Language Models}

\begin{abstract}
Mathematical reasoning remains a particularly formidable challenge for language models, largely due to its intricate and highly structured requirements. This work presents DeepSeekMath~7B, an extension of DeepSeek-Coder-Base-v1.5 7B that is further pre-trained with a specialized dataset comprising 120B math-specific tokens, curated from Common Crawl in addition to text and code data. DeepSeekMath~7B attains a remarkable 51.7\% on the challenging MATH competition benchmark, approaching the capabilities of Gemini-Ultra and GPT-4, and does so without leveraging external tools or ensembling strategies. Using self-consistency over 64 samples, DeepSeekMath~7B reaches 60.9\% on MATH. The notable proficiency of DeepSeekMath~in mathematical reasoning is primarily due to two critical innovations: (1) effective utilization of vast public web resources through a carefully designed data selection mechanism, and (2) the development of Group Relative Policy Optimization (GRPO), a novel adaptation of Proximal Policy Optimization (PPO), which simultaneously bolsters mathematical reasoning and makes PPO more memory efficient.
\end{abstract}

\section{Introduction}

Large language models (LLM) have transformed the landscape of artificial intelligence for mathematical reasoning, leading to substantial improvements on quantitative reasoning benchmarks \citep{MATH} and geometry-based benchmarks \citep{trinh2024solving}. Additionally, such models have become essential tools for aiding humans in addressing complex mathematical tasks \citep{tao}. Nonetheless, premier models like GPT-4 \citep{gpt4} and Gemini-Ultra \citep{gemini} remain closed-source, while the current suite of publicly available open-source models still exhibits a significant gap in performance.

This paper presents DeepSeekMath, a specialized language model tailored for mathematics that markedly surpasses other open-source models in mathematical reasoning and approaches GPT-4 performance across academic benchmarks. Central to our approach is the construction of the DeepSeekMath Corpus, a substantial, high-quality pre-training dataset containing 120B math tokens. This corpus is curated from Common Crawl (CC) data by employing a classifier based on fastText \citep{joulin2016fasttext}. For the first training round, we use positive samples from OpenWebMath \citep{openwebmath} and negative samples from a varied assortment of other web pages to train the classifier. The trained classifier is then deployed to identify additional relevant examples from CC, which are subsequently verified and refined by human annotators. With this updated, richer dataset, the classifier undergoes further training to enhance its accuracy. Evaluation demonstrates the corpus’s quality: our foundational model, DeepSeekMath-Base 7B, obtains a 64.2\% score on GSM8K \citep{gsm8k} and 36.2\% on the competition-grade MATH dataset \citep{MATH}, outperforming Minerva 540B \citep{minerva}. Additionally, since the DeepSeekMath Corpus is multilingual, improvements are observed on Chinese mathematical benchmarks \citep{wei2023cmath,agieval}. We regard our contributions to mathematical data curation as an initial framework for the broader research field, acknowledging that considerable advances remain possible.

DeepSeekMath-Base is constructed by initializing from DeepSeek-Coder-Base-v1.5 7B \citep{deepseek-coder}, as empirical evidence suggests that leveraging a code-focused pretrained model is preferable to starting with a general-purpose LLM. In addition, our results show that training on mathematical data enhances performance not only in math-specific tasks but also on broader benchmarks such as MMLU \citep{mmlu} and BBH \citep{bbh}. This indicates that math training extends beyond augmenting mathematical competencies, yielding broader gains in general reasoning skills.

Upon completion of pre-training, we subject DeepSeekMath-Base to mathematical instruction tuning, incorporating chain-of-thought \citep{cot}, program-of-thought \citep{pot,pal}, and tool-integrated reasoning \citep{tora} data. The model resulting from this process, DeepSeekMath-Instruct 7B, surpasses all other 7B-level models and approaches the performance seen in leading 70B open-source instruction-tuned systems.

Moreover, we propose Group Relative Policy Optimization (GRPO), a novel reinforcement learning (RL) strategy derived from Proximal Policy Optimization (PPO) \citep{schulman2017proximal}. Unlike PPO, GRPO dispenses with the critic network and instead relies on group-based baseline estimates, markedly decreasing computational demands. By employing only a subset of English instruction tuning data, GRPO achieves notable improvements over the already robust DeepSeekMath-Instruct, delivering considerable gains in both in-domain (GSM8K: 82.9\% $\rightarrow$ 88.2\%, MATH: 46.8\% $\rightarrow$ 51.7\%) and out-of-domain mathematics tasks (for example, CMATH: 84.6\% $\rightarrow$ 88.8\%) during the RL stage. In addition, we introduce a unified framework for interpreting various methodologies, such as Rejection Sampling Fine-Tuning (RFT) \citep{yuan2023scaling}, Direct Preference Optimization (DPO) \citep{dpo}, PPO, and GRPO. Under this framework, these techniques can be viewed as variants of direct or simplified RL approaches. We conduct thorough experimentation—exploring aspects such as online versus offline learning, result versus process-based supervision, and single-turn versus iterative RL—to uncover the foundational principles underlying this unified scheme. Finally, we elucidate the mechanisms by which RL enhances the performance of instruction-tuned models and outline future research avenues for advancing RL effectiveness within this comprehensive paradigm.

\subsection{Contributions}
This work presents advances in large-scale math pre-training and investigates the role of reinforcement learning through extensive experiments and analysis.

\textbf{Math Pre-Training at Scale}

\begin{itemize}[topsep=0pt]
    \item We demonstrate that the openly available Common Crawl dataset possesses substantial mathematical content. Through the deployment of a carefully engineered data selection framework, we develop the DeepSeekMath Corpus—an extensive dataset containing 120B tokens sourced from web content curated for mathematical relevance. This dataset is nearly 7 times larger than the mathematical webpages employed by Minerva \citep{minerva}, and surpasses the size of OpenWebMath \citep{openwebmath} by a factor of 9.

    \item Our DeepSeekMath-Base 7B pre-trained model achieves a level of performance on par with Minerva 540B \citep{minerva}, suggesting that mathematical reasoning capacity is not exclusively dictated by model size. In fact, pre-training a more compact model on well-curated, high-quality data can yield robust results.

    \item We provide empirical insights from our math training studies. Conducting code pre-training prior to math pre-training measurably boosts the model’s mathematical problem-solving proficiency, with benefits observed both in scenarios involving external tools and those without. This partially addresses the debated issue: \textit{does code training enhance reasoning skills?} Our findings support an affirmative response, at least within the context of mathematical reasoning.

    \item Contrary to common practice in the domain—where arXiv papers are frequently included for training—we observe no meaningful gains across all evaluated mathematical benchmarks by doing so in our experiments.
\end{itemize}

\textbf{Exploration and Analysis of Reinforcement Learning}

\begin{itemize}[topsep=0pt]
    \item We present Group Relative Policy Optimization (GRPO), a reinforcement learning approach designed for both efficacy and computational efficiency. Distinct from traditional algorithms like Proximal Policy Optimization (PPO), GRPO eliminates the need for a critic model by employing group-based score estimations for the baseline, thereby greatly lowering the computational burden required for training.
    \item Our results show that incorporating GRPO into instruction tuning markedly improves the capabilities of DeepSeekMath-Instruct using only instruction-tuning datasets. Additionally, we report notable improvements in generalization to out-of-domain tasks throughout the reinforcement learning phase.
    \item We establish a unified conceptual framework connecting methods such as RFT, DPO, PPO, and GRPO. Extensive experimental analysis—spanning comparisons of online versus offline learning, outcome-based versus process-based supervision, and single-turn versus iterative reinforcement learning—allows us to systematically analyze the critical factors underpinning this paradigm.
    \item Leveraging this unified framework, we delve into the underlying causes that make reinforcement learning effective and articulate several promising pathways for future progress in optimizing reinforcement learning for large language models.
\end{itemize}

\subsection{Summary of Evaluations and Metrics}

\begin{itemize}[topsep=0pt]
    \item \textbf{English and Chinese Mathematical Reasoning}: We thoroughly evaluate our models on both English and Chinese mathematical benchmarks, which encompass problems ranging from elementary to collegiate difficulty. The English suite includes GSM8K \citep{gsm8k}, MATH \citep{MATH}, SAT \citep{llemma}, OCW Courses \citep{minerva}, and MMLU-STEM \citep{mmlu}. The Chinese benchmarks are composed of MGSM-zh \citep{mgsm}, CMATH \citep{wei2023cmath}, Gaokao-MathCloze \citep{agieval}, and Gaokao-MathQA \citep{agieval}. Assessments are performed on both the models' proficiency at generating comprehensive, self-sufficient solutions and on their ability to execute solutions with Python programming.

    On English benchmarks, DeepSeekMath-Base achieves results that rival those of Minerva 540B \citep{minerva}, a proprietary system, and outperforms all current open-source base models (such as Mistral 7B \citep{mistral} and Llemma-34B \citep{llemma})—regardless of their prior math-specific training, and frequently by substantial margins. Importantly, DeepSeekMath-Base demonstrates leading performance on Chinese mathematical tasks, a result attributable to our inclusive pre-training strategy that integrates both English and non-English high-quality data, unlike approaches taken in earlier studies \citep{minerva,llemma} that restrict themselves to English-only data. Following instruction tuning and reinforcement learning, our DeepSeekMath-Instruct and DeepSeekMath-RL models exhibit outstanding capability, setting a new milestone in open-source research by surpassing 50\% accuracy on the highly challenging, competition-grade MATH dataset.
\end{itemize}

\item \textbf{Formal Mathematics}: We assess DeepSeekMath-Base's capability using the informal-to-formal theorem proving benchmark introduced by \citep{dsp_proof}, evaluated on the miniF2F dataset \citep{minif2f} with Isabelle \citep{isabelle} serving as the chosen proof assistant. DeepSeekMath-Base achieves robust few-shot performance in automatic formalization tasks.

\item \textbf{Natural Language Understanding, Reasoning, and Code}: In order to obtain a thorough evaluation of the models' proficiencies in language comprehension, logical reasoning, and coding, DeepSeekMath-Base is benchmarked on the Massive Multitask Language Understanding (MMLU) suite \citep{mmlu}, which comprises 57 diverse multiple-choice problems, as well as BIG-Bench Hard (BBH) \citep{bbh}, featuring 23 problems that primarily demand complex, multi-step reasoning. Additionally, HumanEval \citep{codex} and MBPP \citep{mbpp}, standard assessments for programming language models, are also included. Pre-training with mathematical data contributes positively to both reasoning and language comprehension capabilities.

\section{Math Pre-Training}

\subsection{Data Collection and Decontamination} This section describes how the DeepSeekMath Corpus was compiled from the Common Crawl dataset. Figure \ref{fig:data_collect} illustrates our stepwise methodology, where an iterative pipeline facilitates the large-scale extraction of mathematical content from Common Crawl, beginning with an initial, high-quality seed corpus composed of mathematical datasets. This general framework is equally adaptable to other specialized domains, such as programming.

As an initial resource, we select OpenWebMath \citep{openwebmath}, a high-quality repository of mathematics-related web texts, to seed our corpus. With this dataset, we train a fastText classifier \citep{joulin2016fasttext} to identify additional web pages with characteristics similar to OpenWebMath. Specifically, we sample 500,000 examples from the seed set as positive cases and another 500,000 randomly chosen web pages from Common Crawl to serve as negatives. Using the open-source fastText library\footnote{\url{https://fasttext.cc}}, we set the vector dimensionality to 256, the learning rate at 0.1, a word n-gram window of up to 3, a minimum word frequency threshold of 3, and train for 3 epochs. To curtail the size of the Common Crawl dataset, both URL deduplication and near-duplicate removal methods are applied, producing a pool of 40B HTML documents. The trained fastText model is subsequently used to retrieve web pages likely to contain mathematical content from this refined dataset. To further ensure content quality, the retrieved pages are ranked by their predicted scores, with only the highest scoring entries retained. The extent of the retained data is evaluated through pre-training using token subsets of the top 40B, 80B, 120B, and 160B tokens. For the initial phase, the selection is limited to the highest scoring 40B tokens.

Following the initial round of data collection, a significant number of mathematical web pages remain uncaptured. This gap is largely attributable to the limited diversity in the positive training examples used for the fastText model. To address this, we broaden our selection of mathematical web sources to enhance the seed corpus and thereby refine the fastText model. Our procedure starts by partitioning the entire Common Crawl dataset into distinct domains, where each domain encompasses all web pages under a single base URL. For each domain, we determine the proportion of pages collected during the first pass. Domains in which more than 10\% of pages have already been collected are labeled as math-related (for example, \url{mathoverflow.net}). 

We then manually review and label the specific URLs within these domains that correspond to mathematical content (such as \url{mathoverflow.net/questions}). Previously uncollected web pages linked to these annotated URLs are incorporated into the seed corpus, enabling us to increase the pool of positive examples and train an upgraded fastText model that can identify even more mathematical data in subsequent iterations. After repeating this data collection process four times, we compile a total of 35.5M mathematical web pages, amounting to 120B tokens. Notably, by the fourth iteration, we observe that 98\% of the data had already been acquired in the preceding iteration, prompting us to terminate the collection process.

To mitigate the risk of contaminating evaluation benchmarks, we apply filtering techniques analogous to those in \cite{deepseek-coder}. Specifically, we remove any web pages containing questions or answers from notable English mathematical benchmarks like GSM8K \citep{gsm8k} and MATH \citep{MATH}, as well as Chinese datasets such as CMATH \citep{wei2023cmath} and AGIEval \citep{agieval}. The filtering method operates as follows: any portion of text containing a 10-gram sequence that exactly matches any segment from the benchmark datasets is excluded from the training set. For benchmark excerpts shorter than 10 grams but at least 3 grams in length, we utilize exact matching to screen and eliminate affected web pages.

\subsection{Validating the Quality of the DeepSeekMath~Corpus}

To evaluate the relative quality of the DeepSeekMath~Corpus, we conduct pre-training experiments that benchmark our dataset against the most recent and widely used math-training corpora:

\begin{itemize}[topsep=0pt]
    \item \textbf{MathPile} \citep{mathpile}: an 8.9B-token corpus derived from various sources such as textbooks, Wikipedia, ProofWiki, CommonCrawl, StackExchange, and arXiv, where the latter contributes over 85\% of the data;
    \item \textbf{OpenWebMath} \citep{openwebmath}: a dataset consisting of 13.6B tokens from CommonCrawl, filtered to retain mathematical material;
    \item \textbf{Proof-Pile-2} \citep{llemma}: includes OpenWebMath, AlgebraicStack (with 10.3B tokens of mathematical code), and arXiv papers (28.0B tokens). For Proof-Pile-2 experiments, we adhere to the arXiv:Web:Code 2:4:1 ratio described in \cite{llemma}.
\end{itemize}

\subsubsection{Training Setting}
\label{sec:quality-policy}
We adapt a general-purpose pre-trained language model with 1.3B parameters, structurally aligned with the DeepSeek LLMs framework \citep{deepseek-llm}, and refer to this variant as DeepSeek-LLM 1.3B. Individual models are trained on each of the mathematical datasets, utilizing a total of 150B tokens per dataset. The HAI-LLM framework \citep{haillm}, known for its efficiency and lightweight nature, supports all training experiments. Mirroring the protocol established for DeepSeek LLMs, optimization employs AdamW \citep{adamW} with $\beta_1=0.9$, $\beta_2=0.95$, and $\mathrm{weight\_decay}=0.1$. The learning rate follows a multi-phase schedule: after a 2,000-step warmup phase, it reaches its maximum, reduces to 31.6\% at the 80\% training mark, and declines further to 10.0\% of the peak by 90\% completion. The peak learning rate is set to 5.3e-4, with batches comprising 4M tokens and a maximum sequence length (context window) of 4K tokens.

\subsubsection{Evaluation Results}

\textbf{DeepSeekMath~Corpus exhibits superior quality, extensive multilingual coverage, and stands out as the largest mathematical corpus to date.}

\begin{itemize}[topsep=0pt]
    \item \textbf{High-quality}:
    We assess downstream capabilities across eight mathematical benchmarks employing few-shot chain-of-thought prompting \cite{cot}. Table \ref{tab:corpora_comparison} reveals that models trained on the DeepSeekMath~Corpus achieve consistently higher performance. Additionally, Figure \ref{fig:corpora_comparisons} demonstrates that, at the 50B token point (equivalent to one complete epoch of Proof-Pile-2), the DeepSeekMath-trained model outperforms the Proof-Pile-2 counterpart, suggesting that the average quality of data in DeepSeekMath is higher.
    \item \textbf{Multilingual}:
    DeepSeekMath~Corpus aggregates mathematical resources from diverse languages, with English and Chinese as the predominant ones. Table \ref{tab:corpora_comparison} indicates that models trained on DeepSeekMath~Corpus exhibit improved mathematical reasoning in both English and Chinese. Conversely, existing corpora, mainly composed of English content, deliver minimal benefits and can even negatively affect performance on Chinese-language tasks.
    \item \textbf{Large-scale}:
    The scale of DeepSeekMath~Corpus substantially exceeds that of previous mathematical datasets. Figure \ref{fig:corpora_comparisons} highlights that DeepSeek-LLM 1.3B, when trained on DeepSeekMath, experiences a sharper learning trajectory and sustained performance gains. In contrast, smaller baseline corpora exhaust their novel data quickly due to repeated cycling during training, causing the resulting model performance to stagnate.
\end{itemize}

\subsection{Training and Evaluating DeepSeekMath-Base 7B}

Here, we present DeepSeekMath-Base 7B, a foundational model specifically designed with robust mathematical reasoning capabilities. The initial weights for our model originate from DeepSeek-Coder-Base-v1.5 7B \citep{deepseek-coder}, and subsequent training was performed over a corpus comprising 500B tokens. The composition of the training dataset is as follows: 56\% derives from the DeepSeekMath Corpus, 4\% from AlgebraicStack, 10\% from arXiv, 20\% from Github code repositories, while the remaining 10\% includes English and Chinese natural language data from Common Crawl. The majority of the training configuration mirrors the protocol outlined in Section \ref{sec:quality-policy}; notable deviations include an elevated learning rate, capped at 4.2e-4, and an increased batch size of 10 million tokens.

A thorough evaluation is performed to analyze the mathematical problem-solving skills of DeepSeekMath-Base 7B. We examine the model’s competency in generating independent mathematical solutions without tool assistance, resolving problems with tool integration, and executing formal theorem-proving tasks. In addition to mathematics-specific tests, we broaden our analysis by measuring the model’s capabilities in natural language comprehension, logical reasoning, and programming.

\paragraph{Mathematical Problem Solving with Step-by-Step Reasoning}

The model’s ability to solve mathematical questions is assessed via few-shot chain-of-thought prompting \citep{cot} across eight separate English and Chinese evaluation datasets. These benchmarks address a wide array of mathematical challenges: quantitative reasoning tasks such as those found in GSM8K \citep{gsm8k}, MATH \citep{MATH}, and CMATH \citep{wei2023cmath}; as well as multiple-choice assessments like MMLU-STEM \citep{mmlu} and Gaokao-MathQA \citep{agieval}. Together, these datasets span a spectrum from elementary to collegiate-level mathematics, encompassing a broad range of domains within the discipline.

As presented in Table \ref{tab:base_math_cot}, DeepSeekMath-Base 7B attains top results across all eight evaluated benchmarks within the cohort of open-source base models (comprising the popular general-purpose Mistral 7B \citep{mistral} and the more recent Llemma 34B \citep{llemma}, which was math-trained using the Proof-Pile-2 corpus \citep{llemma}). In particular, DeepSeekMath-Base achieves a margin exceeding 10\% absolute accuracy over previous open-source base models on the challenging competition-grade MATH dataset and surpasses Minerva 540B \citep{minerva}, a closed-source base model with 77 times more parameters, which is constructed on PaLM \citep{palm} and further optimized on mathematical literature.

\paragraph{Mathematical Problem Solving with Tool Use} To assess program-assisted mathematical reasoning, we apply few-shot program-of-thought prompting \citep{pot,pal} to GSM8K and MATH datasets. In this evaluation, models are instructed to answer each problem by generating a Python script, utilizing modules such as \textit{math} and \textit{sympy} for complex calculations. The output from program execution is taken as the solution. According to Table \ref{tab:base_math_pot_proof}, DeepSeekMath-Base 7B establishes a new state-of-the-art among open-source models, outperforming the previous best Llemma 34B.

\paragraph{Formal Mathematics}

Automated formal proof generation is increasingly recognized for its ability to improve the reliability and efficiency of mathematical proof processes. To this end, DeepSeekMath-Base 7B is benchmarked on informal-to-formal proof conversion from \citep{dsp_proof}, a task involving the transformation of an informal statement, its formal equivalent, and a natural language proof into a formal proof. Evaluation is conducted using miniF2F \citep{minif2f}, a formal Olympiad mathematics benchmark, by prompting the model to produce a formal proof in Isabelle for each item via few-shot examples. Following the protocol in \cite{dsp_proof}, the model outputs proof sketches, after which Sledgehammer \citep{sledgehammer}, an automated proof assistant, is employed to complete remaining proof steps. Results in Table \ref{tab:base_math_pot_proof} reveal that DeepSeekMath-Base 7B is highly competent in the task of autoformalization.

\paragraph{Natural Language Understanding, Reasoning, and Code}

To gauge the model's abilities in natural language comprehension, reasoning, and programming, we evaluate performance on MMLU \citep{mmlu}, BBH \citep{bbh}, HumanEval \citep{codex}, and MBPP \citep{mbpp} respectively. Table \ref{tab:understanding_reasoning_code} shows that DeepSeekMath-Base 7B achieves pronounced gains over its predecessor DeepSeek-Coder-Base-v1.5 \citep{deepseek-coder} on MMLU and BBH, underscoring the efficacy of mathematical pretraining in bolstering language understanding and logical reasoning. Furthermore, due to the incorporation of code-specific tokens during continual training, DeepSeekMath-Base 7B preserves the code generation capabilities of DeepSeek-Coder-Base-v1.5 across both code-related benchmarks. Collectively, DeepSeekMath-Base 7B markedly exceeds the general model Mistral 7B \citep{mistral} on all three tasks related to reasoning and code generation.

\section{Supervised Fine-Tuning}

\subsection{SFT Data Curation}

We assemble an instruction-tuning dataset for mathematics that encompasses both English and Chinese problems, representing a broad array of mathematical domains and varying degrees of difficulty. Each problem is coupled with solutions formulated using chain-of-thought (CoT) \citep{cot}, program-of-thought (PoT) \citep{pot,pal}, or tool-integrated reasoning methods \citep{tora}. In total, our training set comprises 776K examples.

\begin{itemize}[topsep=0pt]
    \item \textbf{English mathematical datasets}: Our English corpus features annotated versions of GSM8K and MATH datasets with tool-integrated solution strategies. In addition, we utilize a selection from MathInstruct \citep{MathInstruct} and the training split of Lila-OOD \citep{lila}, each containing solutions via CoT or PoT approaches. The included problems span a wide variety of mathematical domains, including algebra, probability, number theory, calculus, and geometry.
    \item \textbf{Chinese mathematical datasets}: We collect K-12 mathematics questions in Chinese, covering 76 sub-disciplines such as linear equations. Each item is provided with both CoT-based and tool-integrated reasoning solutions.
\end{itemize}

\subsection{Training and Evaluating DeepSeekMath-Instruct 7B}

This section presents DeepSeekMath-Instruct 7B, a model derived through mathematical instruction tuning from DeepSeekMath-Base. For training, we concatenate samples randomly up to a maximum context window of 4K tokens. Model optimization is carried out over 500 steps using a batch size of 256 and a fixed learning rate of 5e-5.

The model’s quantitative reasoning capabilities, with and without external tool usage, are tested on four benchmarking datasets in both English and Chinese. Our evaluation includes a comparison with top-performing models in the field:

\begin{itemize}[topsep=0pt]
    \item \textbf{Closed-source models} considered in our comparison include: (1) the GPT family—especially GPT-4 \citep{gpt4} and GPT-4 Code Interpreter~\footnote{\url{https://openai.com/blog/chatgpt-plugins\#\#code-interpreter}} for their advanced abilities, (2) Gemini Ultra and Pro \citep{gemini}, (3) Inflection-2 \citep{inflection-2}, (4) Grok-1~\footnote{\url{https://x.ai/model-card}}, in addition to recent entries from Chinese developers such as (5) Baichuan-3~\footnote{\url{https://www.baichuan-ai.com}} and (6) the latest GLM-4~\footnote{\url{https://open.bigmodel.cn/dev/api\#glm-4}} from the GLM suite \citep{glm}.
\end{itemize}

These models are designed for broad applicability and have typically passed through multiple stages of alignment. 
    \item \textbf{Open-source models} featured in our study include: standard generalist LLMs such as (1) DeepSeek-LLM-Chat 67B \citep{deepseek-llm}, (2) Qwen 72B \citep{qwen}, (3) SeaLLM-v2 7B \citep{seallm}, and (4) ChatGLM3 6B \citep{chatglm3}; alongside models with particular enhancements for mathematical capabilities, such as (5) InternLM2-Math 20B~\footnote{\url{https://github.com/InternLM/InternLM-Math}}, which builds on InternLM2 via mathematics-specific pre-training and subsequent instruction tuning, (6) Math-Shepherd-Mistral 7B, which leverages PPO training \citep{schulman2017proximal} on Mistral 7B \citep{mistral} using a process-supervised reward model, (7) the WizardMath suite \citep{wizardmath}, which bolsters the mathematical reasoning skills of Mistral 7B and Llama-2 70B \citep{llama2} through the application of evolve-instruct (an instruction tuning method driven by AI-generated instructions) and PPO training, largely utilizing training data from GSM8K and MATH, (8) MetaMath 70B \citep{metamath}, which is a version of Llama-2 70B fine-tuned on a data-augmented blend of GSM8K and MATH, (9) ToRA 34B \cite{tora}, consisting of CodeLlama 34B adapted for mathematical reasoning with tool integration, and (10) MAmmoTH 70B \citep{MathInstruct}, a Llama-2 70B variant instruction-tuned on the MathInstruct dataset.
\end{itemize}

According to Table \ref{tab:sft_rl_math}, when tool assistance is not permitted during evaluation, DeepSeekMath-Instruct 7B achieves robust results in stepwise reasoning. Particularly on the MATH benchmark, which features competition-level questions, our model outperforms every open-source counterpart and most commercial systems (including Inflection-2 and Gemini Pro) by a margin of at least 9 percentage points. This superiority extends to much larger models (e.g., Qwen 72B) as well as those explicitly reinforced on mathematical problems (such as WizardMath-v1.1 7B). DeepSeekMath-Instruct attains performance levels competitive with Chinese proprietary offerings like GLM-4 and Baichuan-3 on MATH, though it still lags behind leading commercial models such as GPT-4 and Gemini Ultra.

When model evaluation includes both natural language reasoning and tool-assisted programmatic problem solving, DeepSeekMath-Instruct 7B approaches 60\% accuracy on the MATH dataset, establishing a new peak among open-source options. Across other datasets, our system matches or surpasses the performance of DeepSeek-LLM-Chat 67B, the former state-of-the-art open model that is approximately ten times larger.

\section{Reinforcement Learning}

\subsection{Group Relative Policy Optimization} 
Reinforcement learning (RL) has repeatedly demonstrated its potential to further boost the mathematical reasoning of LLMs following Supervised Fine-Tuning (SFT) \citep{wang2023math,wizardmath}. Here, we propose a novel and efficient RL framework termed Group Relative Policy Optimization (GRPO).

\subsubsection{From PPO to GRPO} 
Proximal Policy Optimization (PPO) \citep{schulman2017proximal} stands as a prevalent actor-critic RL approach for policy refinement in LLMs during the RL phase \citep{ouyang2022training}. The method trains LLMs by maximizing the following surrogate loss:

\begin{equation}
\footnotesize
    \mathcal{J}_{PPO}(\theta) = \mathbb{E}{[q \sim P(Q), o \sim \pi_{\theta_{old}}(O|q)]} \frac{1}{|o|} \sum_{t=1}^{|o|} \min \left[ \frac{\pi_\theta(o_{t} | q, o_{<t})}{\pi_{\theta_{old}}(o_{t} | q, o_{<t})} A_{t}, \text{clip} \left( \frac{\pi_\theta(o_{t} | q, o_{<t})}{\pi_{\theta_{old}}(o_{t

Here, $\pi_{\theta}$ and $\pi_{\theta_{old}}$ denote the current and previous policy networks, respectively. The symbols $q$ and $o$ represent questions and outputs, drawn from a dataset of questions and from the historical policy $\pi_{\theta_{old}}$. The constant $\varepsilon$ is a clipping hyperparameter adopted in PPO to enhance training stability. The term $A_t$ refers to the advantage, calculated using Generalized Advantage Estimation (GAE) \citep{gae}, leveraging the reward signals $\{r_{\ge t}\}$ and a trained value estimator $V_{\psi}$. Consequently, PPO necessitates concurrent learning of a value function alongside the policy, and to prevent excessive reward optimization, the typical solution involves applying a per-token KL penalty based on a reference policy within the reward at each generation step, as per \citep{ouyang2022training}:

\begin{equation}
    r_{t} = r_\varphi(q, o_{\le t}) - \beta \log\frac{\pi_{\theta}(o_{t}|q, o_{<t})}{\pi_{ref}(o_{t}|q, o_{<t})},
\label{eq:PPO-reward}
\end{equation}

where $r_\varphi$ designates the reward model, $\pi_{ref}$ stands for the reference policy (commonly the initial SFT model), and $\beta$ quantifies the strength of the KL penalty.

Owing to the fact that the value function in PPO is often an additional neural network of similar scale to the policy itself, this setup incurs considerable overhead in terms of memory and computation. Moreover, in reinforcement learning, the value function is employed as a baseline to reduce variance when computing the advantage. In large language models, however, the reward model is frequently applied to only the final token, thereby complicating accurate value estimation at every step. To tackle these limitations, we introduce, as illustrated in Figure \ref{fig:grpo}, the Group Relative Policy Optimization (GRPO) method. GRPO sidesteps the explicit value function in PPO, instead utilizing the mean reward from multiple outputs—each sampled in response to the same prompt from the old policy—as the baseline. Concretely, for each prompt $q$, GRPO draws a collection $\{o_1, o_2, \cdots, o_G\}$ from $\pi_{\theta_{old}}$ and optimizes the objective below for policy update:

\begin{equation}
\footnotesize
\begin{split}
    \mathcal{J}_{GRPO}(\theta) &= \mathbb{E}{[q \sim P(Q), \{o_i\}_{i=1}^G \sim \pi_{\theta_{old}}(O|q)]}  \\
    & \frac{1}{G}\sum_{i=1}^G\frac{1}{|o_i|} \sum_{t=1}^{|o_i|} \left\{ \min \left[ \frac{\pi_\theta(o_{i,t} | q, o_{i,<t})}{\pi_{\theta_{old}}(o_{i,t} | q, o_{i,<t})} \hat{A}_{i,t}, \text{clip} \left( \frac{\pi_\theta(o_{i,t} | q, o_{i,<t})}{\pi_{\theta_{old}}(o_{i,t} | q, o_{i,<t})}, 1 - \varepsilon, 1 + \varepsilon \right)  \hat{A}_{i,t} \right] - \beta \mathbb{D}_{KL}\left[\pi_{\theta} || \pi_{ref}\right]\right\} ,
\end{split}
\label{eq:GRPO-obj}
\end{equation}

In this formulation, $\varepsilon$ and $\beta$ retain their roles as hyperparameters. The term $\hat{A}_{i,t}$ denotes the advantage calculated using the relative returns from outputs within each group, and its computation is discussed in detail in the subsequent sections. The use of groupwise relative rewards for advantage calculation in GRPO matches the inherently comparative nature of reward modeling, which typically involves datasets that rate responses to the same input. Furthermore, GRPO departs from PPO by adding the KL regularization term directly as a loss, rather than incorporating it into the reward. This direct regularization avoids additional complexity in the computation of $\hat{A}_{i,t}$. Distinct from the KL penalty in (\ref{eq:PPO-reward}), our KL divergence is estimated via the unbiased method outlined in \citep{kl_approx}:

\begin{equation}

\begin{algorithm}[t]
  \small
  \caption{Iterative Group Relative Policy Optimization}
  \textbf{Input} initial policy model $\pi_{\theta_{\text{init}}}$; reward models $r_\varphi$; task prompts $\mathcal{D}$; 
  hyperparameters $\varepsilon$, $\beta$, $\mu$
  \begin{algorithmic}[1]
    \State Initialize policy model: $\pi_\theta \leftarrow \pi_{\theta_{\text{init}}}$
    \For{iteration = 1, \dots, I}
       \State Set the reference model: $\pi_{ref} \leftarrow \pi_{\theta}$
      \For{step = 1, \dots, M}
      \State Draw a batch $\mathcal{D}_b$ from $\mathcal{D}$
      \State Save current policy model: $\pi_{\theta_{old}} \leftarrow \pi_{\theta}$ 
      \State For each prompt $q \in \mathcal{D}_b$, generate $G$ samples $\{o_i\}_{i=1}^G \sim \pi_{\theta_{old}} (\cdot \mid q) $
      \State Evaluate each generated output $o_i$ with $r_{\varphi}$ to obtain rewards $\{r_i\}_{i=1}^G$
      \State For each sampled $o_i$, estimate the group relative advantage $\hat{A}_{i,t}$ at the $t$-th token position.
      \For{GRPO iteration = 1, \dots, $\mu$}
        \State Adjust $\pi_{\theta}$ by maximizing the GRPO objective in Equation \ref{eq:GC-GRPO}
      \EndFor
    \EndFor 
    \State Continuously update $r_\varphi$ using replay-based training.
    \EndFor 
  \end{algorithmic}
  \textbf{Output} $\pi_\theta$
  \label{alg:iter-grpo}
\end{algorithm}

\subsubsection{Outcome Supervision RL with GRPO} In the outcome supervision setting, for any given question $q$, the old policy model $\pi_{\theta_{old}}$ is used to generate a group of outputs $\{o_1, o_2, \cdots, o_G\}$. Each output is evaluated by the reward model, resulting in the reward set $\mathbf{r} = \{r_1, r_2, \cdots, r_G\}$. To normalize these, the group mean is subtracted from each reward and the result is divided by the group standard deviation. Under outcome supervision, the normalized reward for each output $o_i$ is supplied at the output's final token; this same normalized value is then assigned as the advantage $\hat{A}_{i,t}$ to every token in $o_i$: $\hat{A}_{i,t} = \widetilde{r}_i = \frac{r_i- {\rm mean}(\mathbf{r})}{{\rm std}(\mathbf{r})}$. The policy parameters are then updated to maximize the objective described in equation (\ref{eq:GRPO-obj}).

\subsubsection{Process Supervision RL with GRPO} Because outcome supervision gives feedback only at sequence ends, its guidance may not suffice for policies required to perform multi-step mathematical reasoning. Extending the strategy in \cite{wang2023math}, we employ process supervision, which delivers intermediate rewards at the conclusion of each reasoning step. More precisely, given a prompt $q$ and $G$ sampled completions $\{o_1, o_2, \cdots, o_G\}$, the process reward model provides stepwise rewards, yielding $\mathbf{R} = \{ \{r_1^{index(1)},\cdots,r_1^{index(K_1)}\}, \cdots,  \{r_G^{index(1)},\cdots,r_G^{index(K_G)}\} \}$, where $index(j)$ denotes the final token position of the $j$-th step, and $K_i$ is the step count in output $o_i$. These rewards are standardized in the same manner, i.e., $\widetilde{r}_i^{index(j)} = \frac{r_i^{index(j)} - {\rm mean(\mathbf{R})}}{{\rm std(\mathbf{R})}}$. The advantage at each token $t$ within $o_i$ is computed as the sum over all normalized rewards from step indices $j$ with $index(j) \ge t$, that is, $\hat{A}_{i, t} = \sum_{index(j)

\subsubsection{Iterative RL with GRPO} As reinforcement learning advances, the previously trained reward model may no longer provide adequate supervision for the evolving policy model. To address this, we investigate an iterative approach to RL using GRPO. As depicted in Algorithm \ref{alg:iter-grpo}, the iterative GRPO process involves periodically regenerating the reward model’s training dataset from fresh samples drawn by the current policy model. The reward model is then continuously updated through a replay mechanism, incorporating $10\%$ of past data. Subsequently, the policy model is set as the new reference and is further trained with this updated reward model.

\subsection{Training and Evaluating DeepSeekMath-RL}

Our RL experiments are conducted using DeepSeekMath-Instruct 7B. For RL training, we utilize chain-of-thought style questions from GSM8K and MATH found in the SFT dataset, totaling approximately 144K samples. Other SFT questions are excluded to focus on the effect of RL on benchmarks that lack training data during RL. The reward model’s training set is constructed following the procedure of \citep{wang2023math}. The initial reward model is trained using DeepSeekMath-Base 7B with a learning rate of $2 \times 10^{-5}$. For the GRPO method, the policy model employs a learning rate of $1 \times 10^{-6}$ and a KL coefficient of $0.04$. For each prompt, $64$ response samples are generated. The maximum sequence length is set at $1024$, with a batch size of $1024$. The policy model receives a single update after each exploration phase. Evaluation of DeepSeekMath-RL 7B is performed on the same benchmarks as DeepSeekMath-Instruct 7B. Within this context, GSM8K and MATH (in chain-of-thought format) are treated as in-domain evaluation, while the remaining benchmarks are considered out-of-domain.

Table \ref{tab:sft_rl_math} presents the results achieved by both open- and closed-source models employing chain-of-thought and tool-integrated reasoning methods across English and Chinese evaluation sets. Our key observations are as follows: 1) DeepSeekMath-RL 7B achieves accuracy rates of 88.2\% on GSM8K and 51.7\% on MATH when utilizing chain-of-thought prompting, outperforming all open-source competitors within the 7B to 70B parameter range, as well as most closed-source counterparts. 2) Notably, DeepSeekMath-RL 7B’s training is confined solely to chain-of-thought-format instruction tuning data from GSM8K and MATH, initiated from DeepSeekMath-Instruct 7B. Despite this limited training set, DeepSeekMath-RL 7B consistently surpasses DeepSeekMath-Instruct 7B in every metric, thereby highlighting the benefits brought by reinforcement learning.

\section{Discussion}
This section outlines insights obtained during our pre-training and reinforcement learning (RL) experiments.

\subsection{Lessons Learned in Pre-Training}

We begin by detailing the experiences and observations from the pre-training phase. Except where noted otherwise, all training configurations correspond to those described in Section \ref{sec:quality-policy}. For the purposes of this discussion, references to the DeepSeekMath Corpus pertain to an 89B-token dataset curated during the second data collection cycle.

\subsubsection{Code Training Benefits Mathematical Reasoning}

There is a commonly held, though largely anecdotal, belief that training models on code can enhance their reasoning capabilities. Here, we provide partial evidence to support this hypothesis specifically in mathematical contexts: code-focused training improves the proficiency of models in mathematical reasoning, regardless of tool usage.

To evaluate the impact of code-based training on mathematical reasoning, we explored two experimental paradigms: two-stage training and one-stage training. The details of the \textbf{Two-Stage Training} setting are as follows:

\begin{itemize}[topsep=0pt]
    \item \textbf{Code Training for 400B Tokens $\rightarrow$ Math Training for 150B Tokens}: The DeepSeek-LLM 1.3B model is first exposed to 400B code tokens, which is then followed by 150B math token training;
    \item \textbf{General Training for 400B Tokens $\rightarrow$ Math Training for 150B Tokens}: To provide a baseline, we substitute code tokens with general tokens (drawn from an extensive general-purpose dataset curated by DeepSeek-AI) during the initial phase. This setup is used to examine whether code tokens confer greater benefits for mathematical reasoning when compared to general tokens.
\end{itemize}

\textbf{One-Stage Training}

\begin{itemize}[topsep=0pt]
    \item \textbf{Math Training for 150B Tokens}: Here, DeepSeek-LLM 1.3B undergoes training solely on 150B math tokens;
    \item \textbf{Training on a mixture of 400B Code Tokens and 150B Math Tokens}: Since following code pretraining with math fine-tuning has been observed to deteriorate coding performance, we further study the impact of a single-stage approach in which code and math tokens are interleaved during the full training run. The aim is to determine whether such co-training supports mathematical reasoning while simultaneously mitigating catastrophic forgetting.
\end{itemize}

\paragraph{Results}
Tables \ref{tab:code-math} and \ref{tab:code-math-general-eval} present evaluation metrics for downstream tasks across different training protocols.

Pretraining on code data leads to notable gains in program-assisted mathematical reasoning, effective in both the staged and combined training regimens. Specifically, Table \ref{tab:code-math} demonstrates that code token pretraining already offers a substantial boost for GSM8K and MATH tasks involving Python. Subsequent math token training provides additional benefits. Notably, one-stage training with a blended set of code and math tokens largely overcomes catastrophic forgetting observed in two-stage training, while also enhancing both coding ability (Table \ref{tab:code-math-general-eval}) and program-enabled mathematical problem solving (Table \ref{tab:code-math}).

Furthermore, code pretraining is advantageous even for mathematical reasoning tasks that do not utilize tool assistance. The initial code stage in the two-phase protocol generates moderate performance gains and further accelerates improvements during math token training, achieving superior overall results. However, amalgamating code and math tokens within a single-phase training pipeline negatively affects tool-free mathematical reasoning. This could be attributed to the restricted capacity of the 1.3B parameter DeepSeek-LLM model, which may be insufficient to jointly internalize large volumes of both code and mathematics data.

\subsubsection{ArXiv Papers Appear Ineffectual for Enhancing Mathematical Reasoning}

The use of ArXiv papers as a significant portion of mathematical pre-training datasets is widespread \citep{minerva,gpt-f,llemma,mathpile}. Yet, there is a lack of comprehensive studies examining their concrete effect on mathematical reasoning capabilities. Somewhat unexpectedly, our empirical results suggest that arXiv content provides little to no benefit for advancing mathematical reasoning skills. We conducted systematic evaluations using models of varying capacities, notably DeepSeek-LLM 1.3B and DeepSeek-Coder-Base-v1.5 7B \citep{deepseek-coder}, each trained on different arXiv datasets with distinct preprocessing strategies:

\begin{itemize}[topsep=0pt]
    \item \textbf{MathPile} \citep{mathpile}: a corpus of 8.9B tokens, curated using extensive heuristics for data cleaning and filtering, with over 85\% consisting of scientific arXiv articles;
    \item \textbf{ArXiv-RedPajama} \citep{redpajama}: this dataset includes all arXiv LaTeX files after eliminating preambles, comments, custom macros, and bibliographic entries, resulting in a total of 28.0B tokens.
\end{itemize}

In our setup, DeepSeek-LLM 1.3B was trained for 150B tokens and DeepSeek-Coder-Base-v1.5 7B for 40B tokens on each of these corpora. Across the board, models exposed exclusively to arXiv-derived data did not show marked improvements and in some cases even performed worse across a suite of mathematical evaluation tasks with varying degrees of complexity. The evaluation encompassed diverse benchmarks such as GSM8K and MATH for quantitative reasoning (see Table \ref{tab:arxiv-cot}), MMLU-STEM for multiple-choice reasoning (see Table \ref{tab:arxiv-cot}), as well as miniF2F targeting formal mathematical abilities (see Table \ref{tab:arxiv_atp}).

Nonetheless, these findings are preliminary and come with notable caveats. In particular, this work does not investigate:

\begin{itemize}[topsep=0pt]
    \item The possible advantages arXiv data might provide for math-related problems not considered here, such as transforming formal mathematical statements or proofs into their informal counterparts;
    \item The effect of arXiv data when integrated with other data types;
    \item Potential gains that may emerge at larger model sizes beyond those explored in this study.
\end{itemize}

Accordingly, these unresolved questions highlight avenues for future research.

\subsection{Reinforcement Learning Insights}
\subsubsection{Towards a Unified Framework} Here, we articulate a unified framework to systematize and examine distinct training strategies, including SFT, RFT, DPO, PPO, and GRPO. Further, we empirically investigate the determinants influencing this unified approach. Generally, the gradient with respect to the model parameters $\theta$ for a given training procedure can be formally expressed as:

\begin{equation}
    \nabla_{\theta}\mathcal{J}_{\textcolor{red}{\mathcal{A}}}(\theta) = \mathbb{E}[\underbrace{(q,o) \sim \textcolor{red}{\mathcal{D}}}_{Data \ Source}]\left( \frac{1}{|o|} \sum_{t=1}^{|o|}  \underbrace{GC_{{\mathcal{A}}}(q, o, t, \textcolor{red}{\pi_{{rf}}})}_{Gradient \ Coefficient}  \nabla_{\theta}\log \pi_{\theta}(o_t | q, o_{<t})\right).
\label{eq:objective}
\end{equation}

This formulation encompasses three principal elements: 1) \textit{Data Source $\mathcal{D}$}, which supplies the examples for training; 2) \textit{Reward Function $\pi_{{rf}}$}, serving as the generator of training reward signals; and 3) \textit{Algorithm $\mathcal{A}$}, which integrates the data and reward to generate the gradient coefficient $GC$, specifying the degree to which each data point is penalized or rewarded during optimization. The following methods are examined under this general framework:

\begin{itemize}
    \item  \textbf{Supervised Fine-tuning (SFT)}: This approach adapts a pretrained model by training it further on a dataset curated by human annotators.
    \item \textbf{Rejection Sampling Fine-tuning (RFT)}: RFT builds on the SFT model, further refining it by selecting and training on those SFT-generated outputs which are deemed correct, based on pre-established evaluation criteria.
    \item \textbf{Direct Preference Optimization (DPO)}: DPO enhances the SFT model by utilizing pairwise preference comparisons between generated outputs, optimizing the model according to the DPO loss computed on outputs augmented from the SFT model.
    \item \textbf{Online Rejection Sampling Fine-tuning (Online RFT)}: In contrast to standard RFT, this variant initializes from the SFT model but continually updates the policy by sampling augmented responses directly from the current version of the policy model during fine-tuning.
    \item \textbf{PPO/GRPO}: With PPO/GRPO, the policy is initialized with the SFT model and then improved by sampling outputs from the evolving policy and reinforcing learning using these outputs.
\end{itemize}

The elements constituting these approaches are compiled in Table \ref{tab:data_policy}. A comprehensive explanation of the derivation procedures can be found in Appendix \ref{app:analysis-rl}.

\paragraph{Insight into Data Source} We categorize data sources as either online or offline sampling. In the case of online sampling, training data originates from the outputs produced by the currently updated policy during exploration. Offline sampling, by contrast, utilizes training samples generated by the initial SFT model. RFT and DPO are representatives of the offline sampling approach, whereas Online RFT and GRPO utilize online sampling. 

From the comparison illustrated in Figure \ref{fig:rl-analysis}, it is evident that Online RFT yields markedly superior results compared to RFT across two evaluation benchmarks. During the early training epochs, Online RFT's performance is similar to that of RFT, but as training progresses, Online RFT establishes a clear lead, underscoring the effectiveness of the online paradigm. This trend is readily interpretable: early on, there is little deviation between the actor and the SFT model, making their sampled outputs almost indistinguishable. As training advances, however, the data generated by the actor diverges more, and the benefits of sampling from the evolving policy model become more pronounced.

\paragraph{Insight into Gradient Coefficient} In these algorithms, the reward signal is mapped to a gradient coefficient to facilitate parameter updates. For the reward function, our experiments distinguish between `Rule'-based and `Model'-based strategies. The `Rule' method relies on a direct evaluation of answer correctness, whereas the `Model' approach leverages a reward model trained to score responses, with this model’s data sourced from rule-based judgments. Equations \ref{eq:GC-RFT} and \ref{eq:GC-GRPO} reveal an essential distinction: GRPO adaptively modulates its gradient coefficient in response to the reward model’s value, enabling it to reward or penalize responses according to their evaluated merit. Conversely, Online RFT does not implement such dynamic weighting—incorrect responses are not penalized, and all correct responses receive equal reinforcement regardless of score magnitude.

As illustrated in Figure \ref{fig:rl-analysis}, GRPO consistently outperforms online RFT, underscoring the effectiveness of modulating both positive and negative gradient coefficients. Moreover, GRPO+PS achieves better results than GRPO+OS, emphasizing the advantages of adopting finely-grained, step-aware gradient coefficient adjustments. In addition, we investigate the impact of iterative RL, executing two iterative cycles in our experiments. As presented in Figure \ref{fig:iter-rl}, we observe that iterative RL produces substantial performance gains, particularly evident in the first round of iteration.

\subsubsection{Why RL Works?} In our study, we apply reinforcement learning on a selected portion of instruction tuning data, which yields notable improvements relative to the instruction-tuned baseline. To further elucidate the underlying mechanisms of reinforcement learning's efficacy, we compare Pass@K and Maj@K accuracies between the Instruct and RL-enhanced models across two benchmarks. As depicted in Figure \ref{fig:maj-pass}, RL results in an uptick in Maj@K while Pass@K remains largely unchanged. This outcome suggests that reinforcement learning fortifies the overall robustness of the output distribution, i.e., \textbf{the observed gains primarily arise from elevating the likelihood of the correct answer appearing in the TopK predictions rather than fundamentally enhancing base capabilities.} Similarly, \citep{wang2023making} reported a \textbf{misalignment problem} in reasoning within SFT models, and demonstrated that alignment strategies tailored to preferences can bolster reasoning abilities \citep{yuan2023rrhf,song2023preference,wang2023making}.

\subsubsection{How to Achieve More Effective RL?}
\label{sec:effec_rl}
Our results affirm that reinforcement learning is particularly effective for mathematical reasoning. We also introduce a unified framework for conceptualizing various well-known training approaches. In this framework, all such approaches can be interpreted as either variants of RL or simplifications thereof. As stated in Equation \ref{eq:objective}, three fundamental components are identified: Data Source, Algorithm, and Reward Function. Below, we outline possible future research trajectories for these elements.

\paragraph{Data Source} The data source constitutes the foundational input for any training approach. Within the RL context, we specifically consider the data source to be the collection of unlabeled questions accompanied by responses sampled from the policy model. For this study, we limited ourselves to questions utilized during the instruction tuning phase and performed basic nucleus sampling for output generation. We believe this methodological choice partially explains why our RL pipeline only enhances Maj@K accuracy. Moving forward, we plan to evaluate our RL framework using out-of-distribution question prompts, incorporating \textbf{advanced sampling (decoding) strategies}, such as tree-search-based techniques \citep{yao2023tree}. In addition, the application of \textbf{efficient inference methods} \citep{xia-etal-2023-speculative,leviathan2023fast,kwon2023efficient,xia2024unlocking}, which govern the exploration effectiveness of policy models, is anticipated to play a pivotal role in future advancements.

\paragraph{Algorithms} Algorithms employ the observed data and corresponding reward feedback to compute the gradient coefficient, which is then used to update the model's parameters. According to Equation \ref{eq:objective}, current approaches predominantly and \textbf{TRUST} the reward function's guidance, increasing or decreasing the conditional probability of a specific token based on its signal. Nevertheless, the reliability of the reward signal cannot be universally guaranteed, particularly in highly complex tasks. For instance, the PRM800K datasets \citep{lightman2023let}, which have undergone extensive annotation by expert annotators, still exhibit a mislabeling rate of around 20\%\footnote{\url{https://github.com/openai/prm800k/issues/12\#issuecomment-1728491852}}. In response to these challenges, we focus on investigating reinforcement learning algorithms that maintain robustness in the presence of noisy reward signals. We anticipate that such \textbf{WEAK-TO-STRONG} \citep{burns2023weak} alignment frameworks have the potential to fundamentally transform learning algorithm design.

\paragraph{Reward Function} The reward function constitutes the primary source of the training feedback in RL and is typically realized as a neural reward model. We identify three pivotal research directions for reward models: 1) \textbf{Enhancing the generalization capability of reward models.} It is essential for reward models to effectively generalize to unseen or out-of-distribution queries and to advanced decoding outputs; if not, reinforcement learning may only maintain the existing distribution of large language models (LLMs) instead of yielding genuine advancements in their capabilities; 2) \textbf{Capturing and representing reward model uncertainty.} Properly quantifying uncertainty could provide a crucial connection between less reliable reward models and weak-to-strong learning approaches; 3) \textbf{Developing efficient, high-quality process reward models} that deliver fine-grained feedback signals to support sophisticated reasoning processes \citep{lightman2023let,wang2023math}.

\section{Conclusion, Limitation, and Future Work} In this work, we introduce DeepSeekMath, an open-source model that surpasses all existing open-source systems on the competition-grade MATH benchmark and nears the capabilities demonstrated by proprietary models. The model initialization begins with DeepSeek-Coder-v1.5 7B, followed by continuous training over 500B tokens, with a notable portion—120B math tokens—collected from Common Crawl. Through detailed ablation experiments, we find that web-sourced data substantially enriches the quality of mathematical training material, whereas arXiv contributes less than anticipated. Additionally, we propose Group Relative Policy Optimization (GRPO), an adaptation of Proximal Policy Optimization (PPO), which enhances mathematical reasoning proficiency with lower memory requirements. Our empirical analysis confirms GRPO’s advantages, even when DeepSeekMath-Instruct 7B has already attained high benchmark scores. Furthermore, we offer a cohesive framework to interpret several approaches and outline possible avenues for advancing reinforcement learning methodologies.

Despite its high achievement on quantitative reasoning tests, DeepSeekMath demonstrates relatively less strength in geometry and formal theorem-proving compared to closed-source models. For example, during our preliminary evaluation, the model struggled with geometry tasks involving triangles and ellipses, likely pointing to a bias in data selection during pre-training and fine-tuning stages. Moreover, the limitations in model scale hinder DeepSeekMath’s performance on few-shot tasks relative to GPT-4. While GPT-4 shows measurable gains with few-shot examples, DeepSeekMath yields similar outcomes under zero-shot and few-shot settings. Future plans include refining the data selection process to curate an even higher quality pre-training corpus, and further investigating promising strategies for optimizing LLM reinforcement learning (see Section \ref{sec:effec_rl}).

\end{document}